\scanpage{355}
\mspage{\textcircled{18}}

donc une situation ``topologique'' tout à fait spéciale à la catégorie des modèles Cat, qui ne semble pas avoir d'analogue dans des catégories telles que Esp ou $\Delta^{\wedge}$, où pourtant les notions d'ouvert et de fermé complémentaire (avec grains de sel \unsure{ds|dans} $\Delta^{\wedge}$) ont un sens.

L'argument \inserted{catégorique} quasi-trivial qui établit que $(**)$ \deleted{induit} est dans $\mathbb W_1$, i.e. induit une équivalence de catégories entre $\mathrm{Rev}(X)$ et $\mathrm{R\acute{e}v}(\unsure{\mathcal X|\mathscr X})\overset{\sim}{\simeq}\mathrm{Rev}(X_0)\mathop{\times}\limits^{2}_{\mathrm{Rev}(\widetilde X_{01})}\mathrm{Rev}(X_1)$, \unsure{me|ne} semble devoir s'étendre aux dimensions supérieures, à coups de $n$-gerbes \inserted{sur $X$, \unsure{..}} -- mais les fondements sur la théorie des $n$-gerbes \deleted{\unsure{de}} manquent, (Et je n'ai pas le courage d'expliciter l'argument pour $n=2$ (où pourtant Giraud fournit ce qu'il faut.) Aussi je conjecture que $\varphi$ est \deleted{tout} au moins dans $\mathbb W_\infty$, i.e. que le diagramme $(*)$ est $W$-cocartésien si $W\unsure{>|\supset}\mathbb W_\infty$ (et satisfait $\mathrm{Loc}(6'\unsure{.})$).
% slanted
\margin{Mais j'ai fait \unsure{quand même} une esquisse de vérification pour $n=2$, qui me convainc \unsure{qu'on a} bien $\varphi\in\mathbb W_2$, \unsure{Et je|et} \unsure{m’en suis|j’ai} acquis \unsure{de la} \unsure{sorte|certitude} qu'on a même $\varphi\in\mathbb W_\infty$.}

% archive sheet number: 355
